\scanpage{40}
\deleted{elle} pour $r$ fixé, \deleted{elle} \inserted{le fait si elle \deleted{\unsure{est}} \unsure{est|n'est} \inserted{\unsure{ou non|bien}} satisfaite par $u_E^r$} dépend des \deleted{le} choix de $u_E$ et $s_0$ ($\sigma_E$ est ici indiffér., à vrai dire) -- donc du ``module'' $\alpha$.

\deleted{\unsure{On s'intéressera à|On a,}} P. ex. si \inserted{$n=\infty$ mais si} $\alpha$ est tel que la repr. de $D_\infty$ se factorise par $D_r$, donc $u_r=\mathrm{id}$, cette condition ne sera pas satisfaite.

Néanmoins on verra, \inserted{comme on a dit,} que $\psi_r(\alpha)$ peut être défini pour \uline{tt} $\alpha$ -- il se pose alors la question de la signification géom de \deleted{$\psi_r(\alpha)$} la repr. de $D_{\deleted{\unsure{\infty|n}}}$ associé à $\psi_r(\alpha)$, dans le cas ``exceptionnel'' [à caractériser en termes de $\alpha,r$ ! \inserted{(c'est subordonné au Pb général, pour $\alpha$ fixé, de savoir quand trois sommets $s_i,s_j,s_k$ sont alignés\ldots)}] où les $s_{ir}=u^{ir}s_0$ sont tous alignés \ldots.
% slanted
\margin{Question à élucider \ldots !}

$\overline{\underline{\mathrm{VII}}}$ \deleted{Polygônes} \inserted{\uline{Systèmes} \deleted{\unsure{de}}} \uline{réguliers} \deleted{\unsure{finis}} \inserted{\uline{de $n$ sommets}}

% vertical rule beside the parenthesis
(i.e. structures polygonales modulo étoilage \ldots)

Soit $S$ un ensemble \inserted{fini, de cardinal $n$}. On appelle \uline{structure quasi-polygonale} \inserted{(combinatoire)} sur $S$ la donnée d'une structure polygonale
% archive sheet number: 39
